\documentclass{tpks}
\doctitle{Требования к ПО}
\docversion{1}
\docdate{13.11.2026}

\begin{document}
\makeheader
\section{Введение}
\subsection{Назначение}
Спецификация требований к ПО (СПТ) определяет наблюдаемое поведение программной части системы автоматизации управления жизненным циклом релиза ПО. Она служит основанием для проектирования и разработки проверок, уточняя системные требования версии 4 от 13.11.2026. Карточки имеют статус «проект» до независимой проверки и утверждения.
\subsection{Область действия}
Версия 1 охватывает каталог, подготовку релизов, последовательное продвижение, доступ к средам, контроль качества, согласования, учёт установок, откат, управление доступом и аудит. Сборку и установку выполняют внешние системы. ПО хранит сведения и принимает решения о допустимости действий. Материалы авиационного стандарта используются как методическая основа курса; сертификационный уровень учебному продукту не назначен.
\subsection{Определения, акронимы и сокращения}
ID --- идентификатор. SYS --- префикс системного требования, SRD --- требования к ПО. Секунда обозначается «с», процент --- «\%». M --- число различных адресатов согласования, N --- требуемое число положительных решений. Определения взяты из общего глоссария.
\printterms{ПО,СПТ,ВУ,Трассируемость,Авторизация,Аутентификация,Программный продукт,Компонент,Версия компонента,Ревизия исходного кода,Дистрибутив,Контрольная сумма,Релиз,Состав релиза,Фиксация состава релиза,Перечень изменений релиза,Маршрут продвижения,Продвижение релиза,Этап,Шаг,Среда развёртывания,Промышленная среда,ПСИ,КТК,Критический дефект,Правило согласования,Согласование,Принцип «четырёх глаз»,Технологическое окно,Установка,Откат релиза,Журнал аудита,Внешняя система сборки,Внешняя система развёртывания,Внешняя система учёта задач,Внешняя служба аутентификации,Внешняя служба доставки сообщений}
Bearer-токен --- передаваемое с запросом подтверждение доступа от внешней службы аутентификации. Попытка шага --- одно выполнение с отдельным результатом; повтор сохраняет прежнюю попытку в истории. Подтверждённый состав получен из достоверного результата установки, отката или проверки инженером эксплуатации. Подтверждённый пустой состав отличается от неизвестного.
\subsection{Ссылки}
\begin{enumerate}
\item Системные требования, версия 4 от 13.11.2026: \path{docs/01_system_requirements.tex}. Разделы 2--3 задают роли, состояния и окружение; разделы 4--5 --- родительские требования.
\item Стандарт на разработку требований к ПО, версия 2 от 13.11.2026.

\path{docs/04_requirements_standard.tex}.

Задаёт карточку требования и формат источника. Пример: \texttt{parent=[SYS_01_01]}.
\item План разработки ПО, версия 3 от 13.11.2026: \path{docs/03_development_plan.tex}. Задаёт подготовку и независимую проверку СПТ.
\item План управления конфигурацией ПО, версия 2 от 13.11.2026: \path{docs/02_scm_plan.tex}. Задаёт идентификацию и изменение редакций.
\item Козырев В. П. Слайды «Требования», с. 18--27, 38, 78.
Файл: \path{materials/kozyrev_slides/Vse_kozyrev.pdf}. Основание правил высокого уровня и трассировки.
\item Рабочие материалы команды: \path{notes/04_dpm_platform_v.md}, часть Б. Поясняют предметную область; при расхождении действуют системные требования версии 4.
\end{enumerate}
\subsection{Обзор}
Раздел 2 описывает границы изделия, пользователей и ограничения. Раздел 3 содержит обязательства ПО с входами, выходами и критериями проверки, включая обратную трассировку. Приложение поясняет совместное действие требований.

\section{Общее описание}
\subsection{Перспективы изделия}
ПО находится между пользователями и внешними системами. Схема окружения дана на рисунке~\ref{fig:context}. Стрелки обозначают обмен, а не внутренние части ПО.
\begin{figure}[ht]
\centering
\setlength{\unitlength}{1mm}
\begin{picture}(160,72)
\put(52,27){\framebox(56,18){\shortstack{ПО управления\\жизненным циклом релиза}}}
\put(53,59){\framebox(54,10){Пользователи}}
\put(72,59){\vector(0,-1){14}}\put(88,45){\vector(0,1){14}}
\put(0,47){\framebox(43,12){\shortstack{Система\\сборки}}}
\put(43,51){\line(1,0){9}}\put(52,51){\vector(0,-1){6}}
\put(0,8){\framebox(43,12){\shortstack{Система\\развёртывания}}}
\put(43,16){\line(1,0){15}}\put(58,16){\vector(0,1){11}}
\put(69,27){\line(0,-1){15}}\put(69,12){\vector(-1,0){26}}
\put(117,47){\framebox(43,12){\shortstack{Служба\\аутентификации}}}
\put(104,45){\line(0,1){6}}\put(104,51){\vector(1,0){13}}
\put(117,8){\framebox(43,12){\shortstack{Система\\учёта задач}}}
\put(99,27){\line(0,-1){15}}\put(99,12){\vector(1,0){18}}
\end{picture}
\caption{Окружение ПО; обмен со службой доставки сообщений в версии 1 отсутствует}\label{fig:context}
\end{figure}
\subsection{Функции изделия}
Каталог связывает продукты, компоненты, версии и среды. Релиз объединяет версии одного продукта и после фиксации имеет неизменный состав. Маршрут задаёт последовательность этапов и шагов. Продвижение проверяет качество, согласования, занятость и время доступа к среде. Учёт установок отражает подтверждённые сведения; откат восстанавливает предусмотренный состав. Доступ ограничен ролями, изменения отражаются в аудите.
\subsection{Характеристики пользователя}
Пользователь знаком с продуктами и релизами своей организации. Восемь ролей и полномочия приведены в таблице~\ref{tab:roles}. Общий просмотр каталога, релизов, результатов и установок доступен всем восьми ролям. Совмещение ролей сохраняет запрет самосогласования.
\begin{tabl}{|P{3.7cm}|L|}{Пользователи и полномочия}[tab:roles]{\textbf{Роль} & \textbf{Полномочия}}
Администратор & Регистрация продуктов, компонентов, сред; назначение ролей. \\ \hline
Релиз-менеджер & Состав, маршрут, КТК и согласования; запуск, повтор, отмена, запрос отката и его время ожидания. \\ \hline
Разработчик & Общий просмотр. \\ \hline
Инженер по тестированию & Результат назначенного ручного шага непромышленного этапа, включая ПСИ. \\ \hline
Инженер эксплуатации & Промышленные окна, назначенные ручные шаги промышленного этапа; подтверждение состава и разрешение работ после ошибки отката. \\ \hline
Деливери-менеджер & Подтверждение промышленного выпуска после проверок до конца окна. \\ \hline
Согласующий & Одно решение по адресованному открытому запросу. \\ \hline
Аудитор & Просмотр журнала аудита. \\ \hline
\end{tabl}
\subsection{Ограничения на проектирование}
Основание --- системные требования версии 4, разделы 3 и 6. Один маршрут продукта и последовательные этапы сокращают число сценариев. Восемь фиксированных ролей ограничивают объём модели доступа. Сборка и установка переданы внешним системам; собственные средства сборки и хранения дистрибутивов не разрабатываются. Принятый способ входа --- Bearer-токен службы Keycloak. Это установленное системное ограничение.

Специальные устройства отсутствуют. Версия 1 работает для одной организации. Параллельные этапы одного релиза, произвольная цель отката, обход КТК, изменение параметров при повторе и повторное согласование после отказа находятся за границами версии 1. Внутренняя архитектура и требования низкого уровня здесь не задаются.
\subsection{Допущения и зависимости}
Система развёртывания сообщает идентичность задания и позволяет запросить его состояние. Потеря связи не равна ошибке установки. При невозможности узнать результат сохраняется блокировка. Для проверки используются имитаторы с управляемыми успехом, ошибкой, задержкой, потерей и повтором сообщений. Восстановление рассматривается при сохранном хранилище. Самостоятельные изменения среды вне заданий ПО автоматически не обнаруживаются. Доставка сообщений является резервным расширением.

\section{Конкретные требования}
Карточка содержит одно обязательство ПО; источник имеет формат стандарта. ID SYS определены в документе 01. Отклонение означает отказ с причиной и сохранение прежних предметных данных. Ожидание означает сохранение незавершённого действия с причиной, без сообщения об успехе. Ограничения ролей таблицы~\ref{tab:roles} действуют для всех пользовательских операций.
\subsection{Требования к внешним интерфейсам}
\subsubsection{Пользовательский интерфейс}
Пользователь передаёт действие, идентификаторы объектов, требуемые значения и токен доступа. Получает результат, актуальные данные либо причину отказа или ожидания. Состав данных определяется карточками по функциям. Расположение элементов представления здесь не задаётся.
\subsubsection{Аппаратные интерфейсы}
Прямого обмена со специальной аппаратурой нет. Системное время используется для окон и ожиданий; моменты сравниваются с учётом часового смещения. Передача по сети рассматривается на внешней границе ПО.
\subsubsection{Программные интерфейсы}
Логический обмен описан в таблице~\ref{tab:exchange}. Ответ сопоставляется исходному запросу и предметным объектам. Отсутствие достоверного результата сохраняет неопределённость, а не подтверждает успех.
\begin{tabl}{|P{3cm}|L|L|}{Логический протокол обмена}[tab:exchange]{\textbf{Участник} & \textbf{Сообщение и ответ} & \textbf{Ошибка и повтор}}
Система сборки & Вход: компонент, номер версии, ссылка на дистрибутив, сумма, ревизия и момент сборки. Выход: подтверждение версии или отказ. & Повтор идентичных сведений подтверждает прежнюю запись; конфликт и неполные данные отклоняются. \\ \hline
Система развёртывания & Выход: идентичность попытки, релиз, этап, шаг, среда, действие и целевой состав. Вход: идентичность задания, выполняется / успех / ошибка, подтверждённый состав, отчёт проверок. & Потеря ответа вызывает запрос состояния того же задания; новое действие ожидает выяснения результата. Ответ прежней попытки не завершает новую. \\ \hline
Система учёта задач & Выход: запрос по версиям состава для изменений либо по релизу для КТК. Вход: задачи и исправленные дефекты либо число открытых критических дефектов. & Недоступность даёт неполный перечень или невыполненный назначенный критерий КТК. Повторное чтение не меняет внешние данные. \\ \hline
Служба аутентификации & Проверка переданного Bearer-токена: подлинность, срок действия и личность. Достоверное подтверждение допускает проверку ролей. & Недостоверное либо отсутствующее подтверждение приводит к отказу операции. \\ \hline
Служба доставки сообщений & Обмена в обязательной версии 1 нет. & Доставка не является условием допуска релиза. \\ \hline
\end{tabl}
Версия в составе обозначена компонентом и номером версии. Состояние «выполняется» отличается от конечного успеха и ошибки. Неполное сообщение не подтверждает установку. Новая попытка шага отличается от прежней.
\subsubsection{Интерфейсы связи}
Логические требования определяют состав данных и реакции независимо от транспорта. В системных требованиях предписано определить защиту соединений и подлинность внешних отправителей; согласованных контрактов пока нет.

\clarify{Системному аналитику и владельцам внешних систем согласовать транспорт, форматы сообщений, защиту соединений и способ проверки отправителя для сборки, развёртывания и учёта задач. Причина: контрактов во входных документах нет. Снять пометку после утверждения контрактов и проверки совместимости имитаторов. До этого интерфейсный раздел остаётся проектом.}

Идентификатор --- непустое значение, однозначное в своей области. Номер версии уникален в компоненте. Момент времени включает дату, время и часовое смещение. Счётчики проверок и дефектов --- целые неотрицательные числа; успешных не больше общего. Порог доли лежит в [0;100] процентов. Ожидание --- положительное число секунд. Длины строк, точность времени, предельные размеры сообщений и сроки ожидания внешнего чтения должны быть согласованы в контрактах; произвольные ограничения не вводятся.
\subsection{Функциональные требования}
Группы соответствуют функциям 01--10. Штатное выполнение, вмешательство, неизвестный результат и восстановление рассматриваются отдельно.
\subsubsection{Каталог и версии}

\begin{req}{SRD_01_01}{1}
По запросу администратора с новым идентификатором ПО должно регистрировать продукт под этим идентификатором, отклоняя повтор.
\reqsource{parent=[SYS_01_01]}
\reqverify{Испытание: Новый продукт P создан; повтор P отклонён.}
\reqnote{Вход: идентификатор продукта и запрос администратора. Выход: новый продукт либо отказ при повторе. Тип: функциональное. Статус: проект.}
\end{req}

\begin{req}{SRD_01_02}{1}
По запросу администратора ПО должно регистрировать компонент с привязкой к существующему продукту, отклоняя неизвестную привязку.
\reqsource{parent=[SYS_01_02]}
\reqverify{Испытание: Компонент C создан для P; неизвестный Q отклонён.}
\reqnote{Вход: компонент и идентификатор продукта. Выход: компонент с привязкой либо отказ. Тип: функциональное. Статус: проект.}
\end{req}

\begin{req}{SRD_01_03}{1}
По запросу администратора ПО должно регистрировать среду с новым идентификатором, отклоняя повтор идентификатора.
\reqsource{parent=[SYS_01_03]}
\reqverify{Испытание: Создать E дважды; сохранена одна E.}
\reqnote{Вход: идентификатор среды и запрос администратора. Выход: новая среда либо отказ при повторе. Тип: функциональное. Статус: проект.}
\end{req}

\begin{req}{SRD_01_04}{1}
По полному сообщению сборки известного компонента ПО должно регистрировать номер версии, ссылку на дистрибутив, контрольную сумму, ревизию и момент сборки без изменения полученных значений.
\reqsource{parent=[SYS_01_04]}
\reqverify{Испытание: Передать все поля; каждое значение каталога совпадает с сообщением.}
\reqnote{Вход: компонент, номер версии, ссылка, сумма, ревизия и время сборки. Выход: зарегистрированная версия со всеми полученными значениями. Тип: функциональное. Статус: проект.}
\end{req}

\begin{req}{SRD_01_05}{1}
При неизвестном компоненте в сообщении сборки ПО должно отклонять регистрацию версии.
\reqsource{parent=[SYS_01_05]}
\reqverify{Испытание: Неизвестный C отклонён, новой версии нет.}
\reqnote{Вход: сообщение сборки с неизвестным компонентом либо пропущенным обязательным значением. Выход: отказ с причиной; каталог версий сохранён. Тип: функциональное. Статус: проект.}
\end{req}

\begin{req}{SRD_01_06}{1}
При отсутствии любого обязательного значения сборки ПО должно отклонять регистрацию с указанием отсутствующего значения.
\reqsource{parent=[SYS_01_05]}
\reqverify{Испытание: Поочерёдно убрать компонент, номер, ссылку, сумму, ревизию и время; каждый запрос отклонён.}
\reqnote{Вход: сообщение сборки с неизвестным компонентом либо пропущенным обязательным значением. Выход: отказ с причиной; каталог версий сохранён. Тип: функциональное. Статус: проект.}
\end{req}

\begin{req}{SRD_01_07}{1}
При повторном получении идентичных сведений о паре «компонент, номер версии» ПО должно подтверждать существующую запись вместо создания новой.
\reqsource{parent=[SYS_01_06]}
\reqverify{Испытание: Два одинаковых сообщения дают одну версию и один идентификатор записи.}
\reqnote{Вход: сведения сборки для существующей пары «компонент, номер версии». Выход: подтверждение прежней записи либо отказ при конфликте. Тип: функциональное. Статус: проект.}
\end{req}

\begin{req}{SRD_01_08}{1}
При отличии сведений о ранее зарегистрированной паре «компонент, номер версии» ПО должно отклонять сообщение с сохранением прежней записи.
\reqsource{parent=[SYS_01_06]}
\reqverify{Испытание: Изменить ссылку, сумму, ревизию и время по отдельности; исходная запись сохранена.}
\reqnote{Вход: сведения сборки для существующей пары «компонент, номер версии». Выход: подтверждение прежней записи либо отказ при конфликте. Тип: функциональное. Статус: проект.}
\end{req}

\begin{req}{SRD_01_09}{1}
По запросу любой из восьми ролей ПО должно возвращать сохранённые продукты, компоненты, версии, среды, состав и состояние релизов.
\reqsource{parent=[SYS_01_07]}
\reqverify{Испытание: Каждая роль видит подготовленные данные; пользователь без роли отказан.}
\reqnote{Вход: запрос просмотра от пользователя с назначенной ролью. Выход: сведения каталога, состава и состояния релизов. Тип: функциональное. Статус: проект.}
\end{req}

\subsubsection{Релизы}

\begin{req}{SRD_02_01}{1}
По запросу релиз-менеджера ПО должно создавать релиз известного продукта в состоянии «Черновик», отклоняя неизвестный продукт.
\reqsource{parent=[SYS_02_01]}
\reqverify{Испытание: Для P создан черновик P; неизвестный Q отклонён.}
\reqnote{Вход: продукт и запрос создания релиза. Выход: идентификатор черновика либо отказ. Тип: функциональное. Статус: проект.}
\end{req}

\begin{req}{SRD_02_02}{1}
По запросу релиз-менеджера ПО должно включать допустимую зарегистрированную версию в состав черновика.
\reqsource{parent=[SYS_02_02]}
\reqverify{Испытание: Запрошенная версия появилась в составе ровно один раз.}
\reqnote{Вход: черновик и добавляемая зарегистрированная версия. Выход: обновлённый состав черновика. Тип: функциональное. Статус: проект.}
\end{req}

\begin{req}{SRD_02_03}{1}
По запросу релиз-менеджера ПО должно удалять указанную включённую версию из черновика с сохранением остальных позиций.
\reqsource{parent=[SYS_02_03]}
\reqverify{Испытание: Из двух компонентов удалить один; второй сохранён.}
\reqnote{Вход: черновик и удаляемая включённая версия. Выход: состав без выбранной позиции. Тип: функциональное. Статус: проект.}
\end{req}

\begin{req}{SRD_02_04}{1}
По запросу фиксации непустого допустимого черновика ПО должно фиксировать его текущий состав с переходом в «Сформирован».
\reqsource{parent=[SYS_02_04]}
\reqverify{Испытание: Обе позиции непустого состава сохранены, состояние «Сформирован».}
\reqnote{Вход: черновик и запрос фиксации. Выход: зафиксированный непустой состав и «Сформирован» либо отказ пустому составу. Тип: функциональное. Статус: проект.}
\end{req}

\begin{req}{SRD_02_05}{1}
При запросе фиксации пустого черновика ПО должно отклонять фиксацию с сохранением состояния «Черновик».
\reqsource{parent=[SYS_02_04]}
\reqverify{Испытание: Пустой состав и состояние остались прежними, получена причина отказа.}
\reqnote{Вход: черновик и запрос фиксации. Выход: зафиксированный непустой состав и «Сформирован» либо отказ пустому составу. Тип: функциональное. Статус: проект.}
\end{req}

\begin{req}{SRD_02_06}{1}
После фиксации состава ПО должно отклонять добавление, удаление и замену его позиций.
\reqsource{parent=[SYS_02_05]}
\reqverify{Испытание: Все три вида изменения отклонены; состав прежний во всех состояниях после фиксации.}
\reqnote{Вход: запрос изменения зафиксированного состава. Выход: отказ; состав сохранён. Тип: функциональное. Статус: проект.}
\end{req}

\begin{req}{SRD_02_07}{1}
При неизвестной версии, компоненте другого продукта или второй версии включённого компонента ПО должно отклонять добавление в состав.
\reqsource{parent=[SYS_02_06]}
\reqverify{Испытание: Каждый из трёх классов входа отказан; прежний состав сохранён.}
\reqnote{Вход: неизвестная версия, чужой компонент либо вторая версия компонента. Выход: отказ; состав сохранён. Тип: функциональное. Статус: проект.}
\end{req}

\begin{req}{SRD_02_08}{1}
По версиям текущего состава ПО должно получать перечень связанных задач и исправленных дефектов без повторов идентификаторов.
\reqsource{parent=[SYS_02_07]}
\reqverify{Испытание: Общая задача двух версий присутствует один раз; посторонние задачи отсутствуют.}
\reqnote{Вход: версии текущего состава и ответы учёта задач. Выход: перечень связанных задач и исправленных дефектов без повторов. Тип: функциональное. Статус: проект.}
\end{req}

\begin{req}{SRD_02_09}{1}
При недоступности данных учёта задач хотя бы для одной версии состава ПО должно обозначать перечень изменений как неполный.
\reqsource{parent=[SYS_02_08]}
\reqverify{Испытание: Потеря ответа одной версии даёт отметку неполноты, включая пустой результат.}
\reqnote{Вход: недоступные сведения хотя бы одной версии. Выход: признак неполного перечня изменений. Тип: функциональное. Статус: проект.}
\end{req}

\subsubsection{Маршрут продвижения}

\begin{req}{SRD_03_01}{1}
По запросу релиз-менеджера ПО должно сохранять порядок этапов маршрута продукта и их зарегистрированные среды, отклоняя неизвестную среду.
\reqsource{parent=[SYS_03_01]}
\reqverify{Испытание: Три этапа читаются в заданном порядке; неизвестная среда отказана.}
\reqnote{Вход: продукт, список этапов и их среды. Выход: упорядоченный маршрут либо отказ неизвестной среде. Тип: функциональное. Статус: проект.}
\end{req}

\begin{req}{SRD_03_02}{1}
По запросу релиз-менеджера ПО должно сохранять автоматический шаг с действием «установка» либо «автоматическая проверка» и положительным ожиданием в секундах, отклоняя другое действие и неположительное ожидание.
\reqsource{parent=[SYS_03_02]}
\reqverify{Испытание: Оба действия с 60 с приняты; неизвестное действие, 0 и минус 1 с отказаны.}
\reqnote{Вход: вид автоматического действия и время ожидания в секундах. Выход: настройка автоматического шага либо отказ. Тип: функциональное. Статус: проект.}
\end{req}

\begin{req}{SRD_03_03}{1}
По запросу релиз-менеджера ПО должно сохранять ручной шаг с инженером по тестированию для непромышленного этапа либо инженером эксплуатации для промышленного, отклоняя неподходящую роль.
\reqsource{parent=[SYS_03_03]}
\reqverify{Испытание: Правильные назначения приняты; обратные и пользователь без нужной роли отказаны.}
\reqnote{Вход: среда ручного шага и назначаемый исполнитель. Выход: назначение исполнителя либо отказ по роли. Тип: функциональное. Статус: проект.}
\end{req}

\begin{req}{SRD_03_04}{1}
По запросу релиз-менеджера ПО должно сохранять заданный порядок шагов этапа.
\reqsource{parent=[SYS_03_04]}
\reqverify{Испытание: Сохранить три шага; при чтении порядок совпадает с заданным.}
\reqnote{Вход: список шагов этапа в заданном порядке. Выход: сохранённый порядок шагов. Тип: функциональное. Статус: проект.}
\end{req}

\begin{req}{SRD_03_05}{1}
При запуске продвижения ПО должно закреплять маршрут, шаги, КТК и правила согласования по состоянию на момент запуска.
\reqsource{parent=[SYS_03_05]}
\reqverify{Испытание: Изменить все четыре настройки продукта; действующий релиз использует прежние, следующий получает новые.}
\reqnote{Вход: настройки маршрута в момент запуска. Выход: закреплённый маршрут продвижения. Тип: функциональное. Статус: проект.}
\end{req}

\begin{req}{SRD_03_06}{1}
При настройке маршрута ПО должно предоставлять обозначение одного непромышленного этапа как ПСИ с обязательной ручной приёмкой назначенным инженером по тестированию.
\reqsource{parent=[SYS_03_06]}
\reqverify{Испытание: Ручной результат исполнителя завершает приёмку; автоматический отчёт её не завершает.}
\reqnote{Вход: этап маршрута, ручная приёмка и инженер по тестированию. Выход: настройка этапа ПСИ. Тип: функциональное. Статус: проект.}
\end{req}

\subsubsection{Продвижение релиза}

\begin{req}{SRD_04_01}{1}
По запросу релиз-менеджера ПО должно запускать сформированный релиз только при непустом маршруте со шагами на каждом этапе, последнем промышленном этапе с установкой, одном ПСИ с ручной приёмкой перед ним, достаточном числе допустимых согласующих и положительном ожидании автоотката.
\reqsource{parent=[SYS_04_01]}
\reqverify{Испытание: Допустимый набор даёт «На продвижении»; поочерёдно нарушить каждое условие и исходное состояние --- отказ с причиной.}
\reqnote{Вход: сформированный релиз, маршрут, допустимые согласующие и ожидание автоотката. Выход: «На продвижении» либо отказ с нарушенным условием. Тип: функциональное. Статус: проект.}
\end{req}

\begin{req}{SRD_04_02}{1}
ПО должно начинать каждый этап после первого только при состоянии предыдущего «Пройден».
\reqsource{parent=[SYS_04_02]}
\reqverify{Испытание: при выполнении или ошибке предыдущего этапа следующий ожидает. После «Пройден» действуют остальные условия допуска.}
\reqnote{Вход: результат предыдущего этапа. Выход: допуск следующего к остальным проверкам либо ожидание. Тип: функциональное. Статус: проект.}
\end{req}

\begin{req}{SRD_04_03}{1}
ПО должно начинать следующий шаг этапа только после успеха предыдущего.
\reqsource{parent=[SYS_04_03]}
\reqverify{Испытание: Задержка и ошибка первого блокируют второй; успех разрешает второй по порядку.}
\reqnote{Вход: результат предыдущего шага. Выход: начало следующего шага либо ожидание. Тип: функциональное. Статус: проект.}
\end{req}

\begin{req}{SRD_04_04}{1}
При начале автоматического шага ПО должно передавать задание с релизом, средой, действием и зафиксированным составом.
\reqsource{parent=[SYS_04_04]}
\reqverify{Испытание: Имитатор получает ровно значения текущего релиза, среды и шага.}
\reqnote{Вход: активный автоматический шаг, релиз, среда и состав. Выход: внешнее задание и его связь с попыткой. Тип: функциональное. Статус: проект.}
\end{req}

\begin{req}{SRD_04_05}{1}
По конечному результату соответствующего задания ПО должно устанавливать результат его попытки шага «Успешно» или «Ошибка» согласно ответу.
\reqsource{parent=[SYS_04_05]}
\reqverify{Испытание: Оба исхода отражены у своего шага; ответ иной попытки текущую не завершает.}
\reqnote{Вход: идентичность задания и его результат. Выход: результат соответствующей попытки либо сохранение её незавершённости. Тип: функциональное. Статус: проект.}
\end{req}

\begin{req}{SRD_04_06}{1}
Для каждой попытки внешнего действия ПО должно сохранять связь задания с релизом, этапом, шагом и попыткой до ожидания результата.
\reqsource{parent=[SYS_04_04,SYS_04_05,SYS_04_10]}
\reqverify{Испытание: Две попытки различимы; после перезапуска запрос состояния адресован исходному заданию.}
\reqnote{Вход: активный автоматический шаг, релиз, среда и состав. Выход: внешнее задание и его связь с попыткой. Тип: функциональное. Статус: проект.}
\end{req}

\begin{req}{SRD_04_07}{1}
При неизвестном результате ПО должно запрашивать состояние исходного задания без повторной выдачи действия до установления результата.
\reqsource{parent=[SYS_04_10]}
\reqverify{Испытание: Потерять ответ и перезапустить; один исходный запуск и запрос его состояния, без повторного запуска.}
\reqnote{Вход: неизвестный результат и идентичность исходного задания. Выход: запрос состояния без повторного запуска действия. Тип: функциональное. Статус: проект.}
\end{req}

\begin{req}{SRD_04_08}{1}
Если ответ нельзя сопоставить заданию или он не содержит конечного результата, ПО должно сохранять незавершённость текущей попытки.
\reqsource{parent=[SYS_04_05,SYS_04_10]}
\reqverify{Испытание: Чужой ответ и «выполняется» не завершают текущий шаг.}
\reqnote{Вход: идентичность задания и его результат. Выход: результат соответствующей попытки либо сохранение её незавершённости. Тип: функциональное. Статус: проект.}
\end{req}

\begin{req}{SRD_04_09}{1}
ПО должно принимать «выполнен» или «не выполнен» только от назначенного исполнителя активного ручного шага с результатом соответственно «Успешно» или «Ошибка».
\reqsource{parent=[SYS_04_06]}
\reqverify{Испытание: Оба значения исполнителя приняты; иной пользователь, неизвестное значение и неактивный шаг отказаны.}
\reqnote{Вход: исполнитель, ручной шаг и переданный результат. Выход: «Успешно», «Ошибка» либо отказ. Тип: функциональное. Статус: проект.}
\end{req}

\begin{req}{SRD_04_10}{1}
После успеха всех шагов и всех назначенных критериев КТК ПО должно отмечать непромышленный этап «Пройден».
\reqsource{parent=[SYS_04_07]}
\reqverify{Испытание: Каждый неуспешный шаг или критерий блокирует прохождение; без КТК достаточно успеха шагов.}
\reqnote{Вход: результаты шагов, назначенной КТК и промышленное подтверждение. Выход: «Пройден» либо сохранение незавершённого этапа. Тип: функциональное. Статус: проект.}
\end{req}

\begin{req}{SRD_04_11}{1}
ПО должно отмечать промышленный этап «Пройден» только при успехе шагов, всех назначенных критериев КТК и допустимом подтверждении выпуска.
\reqsource{parent=[SYS_04_07]}
\reqverify{Испытание: Без подтверждения --- ожидание; подтверждение при ошибочных шагах или КТК отказано; полный допустимый набор завершает этап.}
\reqnote{Вход: результаты шагов, назначенной КТК и промышленное подтверждение. Выход: «Пройден» либо сохранение незавершённого этапа. Тип: функциональное. Статус: проект.}
\end{req}

\begin{req}{SRD_04_12}{1}
При ошибке шага, невыполненной КТК, отклонённом согласовании или достижении времени ожидания автоматического шага ПО должно переводить релиз в «Требует вмешательства».
\reqsource{parent=[SYS_04_08]}
\reqverify{Испытание: Каждое из четырёх событий даёт вмешательство и блокирует следующий шаг; время отсчитано от передачи задания.}
\reqnote{Вход: ошибка шага, отказ КТК/согласования либо истечение ожидания. Выход: «Требует вмешательства» с причиной. Тип: функциональное. Статус: проект.}
\end{req}

\begin{req}{SRD_04_13}{1}
По запросу релиз-менеджера ПО должно повторять выбранный шаг текущего непройденного этапа требующего вмешательства релиза только при известном завершённом прежнем задании, доступной среде, действующем либо ненужном согласовании и отсутствии назначенного или выполняемого отката.
\reqsource{parent=[SYS_04_09]}
\reqverify{Испытание: Допустимый повтор даёт «На продвижении»; неизвестный результат, отказ согласования и откат блокируют повтор; доступ к промышленной установке учитывает окно.}
\reqnote{Вход: выбранный шаг, известный прежний результат, допуск среды и согласование. Выход: новая попытка прежнего шага и история либо блокировка повтора. Тип: функциональное. Статус: проект.}
\end{req}

\begin{req}{SRD_04_14}{1}
При повторе ПО должно использовать прежние параметры шага с сохранением предшествующих результатов и переносом результатов выбранного и следующих шагов только в историю.
\reqsource{parent=[SYS_04_09]}
\reqverify{Испытание: Повторить второй из трёх; первый сохранён, второй и третий выполняются заново, затем КТК вычисляется заново.}
\reqnote{Вход: выбранный шаг, известный прежний результат, допуск среды и согласование. Выход: новая попытка прежнего шага и история либо блокировка повтора. Тип: функциональное. Статус: проект.}
\end{req}

\begin{req}{SRD_04_15}{1}
После результата «Пройден» последнего этапа ПО должно переводить релиз в «Выпущен».
\reqsource{parent=[SYS_04_11]}
\reqverify{Испытание: До прохождения последнего этапа выпуска нет; после прохождения есть.}
\reqnote{Вход: результат «Пройден» последнего этапа. Выход: «Выпущен». Тип: функциональное. Статус: проект.}
\end{req}

\begin{req}{SRD_04_16}{1}
По запросу релиз-менеджера ПО должно отменять релиз из «Черновик», «Сформирован», «На продвижении» или «Требует вмешательства» только без выполняющегося шага, неизвестного задания, назначенного/выполняемого отката и ожидания подтверждения промышленного выпуска.
\reqsource{parent=[SYS_04_12]}
\reqverify{Испытание: Допустимая отмена даёт «Отменён», новые шаги прекращены, учёт установок сохранён; каждое запрещающее условие даёт отказ.}
\reqnote{Вход: запрос отмены, состояние релиза и незавершённые действия. Выход: «Отменён» либо отказ с причиной. Тип: функциональное. Статус: проект.}
\end{req}

\begin{req}{SRD_04_17}{1}
ПО должно отклонять переход состояния вне таблицы состояний системных требований версии 4, раздел 2.4.
\reqsource{parent=[SYS_04_13]}
\reqverify{Испытание: Возврат сформированного в черновик, запуск отменённого и отмена выпущенного отказаны.}
\reqnote{Вход: исходное состояние и запрошенный переход. Выход: отказ недопустимому переходу. Тип: функциональное. Статус: проект.}
\end{req}

\subsubsection{Доступ к средам}

\begin{req}{SRD_05_01}{1}
При конкурирующих запросах ПО должно предоставлять одну среду не более чем одному этапу или операции отката одновременно.
\reqsource{parent=[SYS_05_01]}
\reqverify{Испытание: Два релиза, затем релиз и откат претендуют на одну среду; выполняется одна работа.}
\reqnote{Вход: претенденты на среду, исходные задания, подтверждение состава и выпуска. Выход: единственный владелец, передача права откату либо ожидание. Тип: функциональное. Статус: проект.}
\end{req}

\begin{req}{SRD_05_02}{1}
ПО должно сохранять занятость среды до завершения внешних заданий с известными результатами и подтверждения состава, включая ожидание подтверждения промышленного выпуска.
\reqsource{parent=[SYS_05_01]}
\reqverify{Испытание: Неизвестное задание, неподтверждённый состав и ожидание выпуска по отдельности сохраняют блокировку конкурента.}
\reqnote{Вход: претенденты на среду, исходные задания, подтверждение состава и выпуска. Выход: единственный владелец, передача права откату либо ожидание. Тип: функциональное. Статус: проект.}
\end{req}

\begin{req}{SRD_05_03}{1}
При допуске отката того же релиза ПО должно передавать ему право работы после завершения исходных заданий с известным результатом без допуска постороннего релиза между ними.
\reqsource{parent=[SYS_05_01]}
\reqverify{Испытание: другой релиз ожидает при передаче среды откату. После ошибки допустимо восстановление по сохранённой цели, даже если текущий состав неизвестен.}
\reqnote{Вход: претенденты на среду, исходные задания, подтверждение состава и выпуска. Выход: единственный владелец, передача права откату либо ожидание. Тип: функциональное. Статус: проект.}
\end{req}

\begin{req}{SRD_05_04}{1}
ПО должно оставлять среду свободной при одном лишь ожидании допуска к началу этапа.
\reqsource{parent=[SYS_05_01]}
\reqverify{Испытание: Этап ждёт согласования или окна до начала работ; среда доступна иной допустимой работе.}
\reqnote{Вход: претенденты на среду, исходные задания, подтверждение состава и выпуска. Выход: единственный владелец, передача права откату либо ожидание. Тип: функциональное. Статус: проект.}
\end{req}

\begin{req}{SRD_05_05}{1}
По запросу инженера эксплуатации ПО должно сохранять промышленное окно с датами, временем и часовым смещением при начале раньше окончания, отклоняя равные и обратные границы.
\reqsource{parent=[SYS_05_02]}
\reqverify{Испытание: Интервал через полночь принят; равные и обратные границы отказаны.}
\reqnote{Вход: дата, время и смещение начала и окончания окна. Выход: сохранённый интервал либо отказ. Тип: функциональное. Статус: проект.}
\end{req}

\begin{req}{SRD_05_06}{1}
Для промышленной установки, её повтора и ручного отката ПО должно разрешать начало не раньше начала и строго раньше окончания окна, сохраняя ожидание вне окна.
\reqsource{parent=[SYS_05_03]}
\reqverify{Испытание: Для всех трёх действий проверить до начала, на начале, перед концом, на конце и после; допуск только внутри полуинтервала.}
\reqnote{Вход: вид работы, среда, текущее время и границы окна. Выход: допуск по окну либо ожидание; для непромышленной среды окно не применяется. Тип: функциональное. Статус: проект.}
\end{req}

\begin{req}{SRD_05_07}{1}
Для непромышленной среды ПО должно оценивать допуск независимо от промышленного окна.
\reqsource{parent=[SYS_05_03]}
\reqverify{Испытание: Без окон допустимая непромышленная установка начинается; прочие ограничения сохранены.}
\reqnote{Вход: вид работы, среда, текущее время и границы окна. Выход: допуск по окну либо ожидание; для непромышленной среды окно не применяется. Тип: функциональное. Статус: проект.}
\end{req}

\subsubsection{Контрольные точки качества}

\begin{req}{SRD_06_01}{1}
По запросу релиз-менеджера ПО должно назначать этапу одну КТК с непустым набором из минимальной доли успешных проверок и/или максимума открытых критических дефектов.
\reqsource{parent=[SYS_06_01]}
\reqverify{Испытание: Оба критерия отдельно и вместе приняты; пустой набор отказан.}
\reqnote{Вход: этап и непустой набор критериев. Выход: назначенная КТК либо отказ пустому набору. Тип: функциональное. Статус: проект.}
\end{req}

\begin{req}{SRD_06_02}{1}
ПО должно отклонять порог доли вне [0;100] процентов и порог дефектов, отличный от целого неотрицательного числа.
\reqsource{parent=[SYS_06_02]}
\reqverify{Испытание: 0 и 100 процентов, 0 дефектов приняты; минус 1 и 101 процент, минус 1 и 0,5 дефекта отказаны.}
\reqnote{Вход: процентный порог и/или порог дефектов. Выход: принятые значения либо отказ с причиной. Тип: функциональное. Статус: проект.}
\end{req}

\begin{req}{SRD_06_03}{1}
После шагов текущей попытки ПО должно считать критерий доли выполненным, если сумма успешных проверок по её отчётам, делённая на сумму всех проверок и умноженная на 100, не меньше порога без округления.
\reqsource{parent=[SYS_06_03]}
\reqverify{Испытание: 95 из 100 при пороге 95 проходят, 94 нет; 2 из 3 при 66,67 не проходят; отчёты старых попыток исключены.}
\reqnote{Вход: пороги, отчёты текущей попытки и число открытых критических дефектов. Выход: результат назначенного критерия либо всей КТК согласно карточке. Тип: функциональное. Статус: проект.}
\end{req}

\begin{req}{SRD_06_04}{1}
При проверке КТК ПО должно считать критерий дефектов выполненным при достоверном числе открытых критических дефектов релиза на момент проверки не больше порога.
\reqsource{parent=[SYS_06_03]}
\reqverify{Испытание: При пороге 0 число 0 проходит, 1 нет; закрытые и некритические дефекты исключены.}
\reqnote{Вход: пороги, отчёты текущей попытки и число открытых критических дефектов. Выход: результат назначенного критерия либо всей КТК согласно карточке. Тип: функциональное. Статус: проект.}
\end{req}

\begin{req}{SRD_06_05}{1}
ПО должно считать КТК выполненной только при выполнении всех назначенных критериев.
\reqsource{parent=[SYS_06_03]}
\reqverify{Испытание: При двух назначенных один отказ даёт отказ всей КТК; неназначенный критерий не влияет.}
\reqnote{Вход: пороги, отчёты текущей попытки и число открытых критических дефектов. Выход: результат назначенного критерия либо всей КТК согласно карточке. Тип: функциональное. Статус: проект.}
\end{req}

\begin{req}{SRD_06_06}{1}
ПО должно считать назначенный критерий невыполненным при отсутствии данных, а критерий доли также при нуле всех проверок и противоречивых счётчиках.
\reqsource{parent=[SYS_06_04]}
\reqverify{Испытание: Недоступные данные, отсутствие отчёта, 0 проверок, отрицательный или нецелый счётчик и успешных больше общего дают отказ с причиной.}
\reqnote{Вход: нет данных, число проверок равно нулю либо счётчики противоречивы. Выход: критерий не выполнен, указана причина. Тип: функциональное. Статус: проект.}
\end{req}

\subsubsection{Согласования}

\begin{req}{SRD_07_01}{1}
По запросу релиз-менеджера ПО должно сохранять правило при различных адресатах с ролью согласующего и целом N от 1 до M, отклоняя нарушение этих условий.
\reqsource{parent=[SYS_07_01]}
\reqverify{Испытание: «2 из 3» принято; повторы адресатов, чужая роль, N=0 и N больше M отказаны.}
\reqnote{Вход: список различных адресатов, их роли и целое N. Выход: правило «N из M» либо отказ. Тип: функциональное. Статус: проект.}
\end{req}

\begin{req}{SRD_07_02}{1}
При N положительных решениях разных допустимых адресатов открытого запроса без отрицательного решения ПО должно закрывать запрос как «согласовано».
\reqsource{parent=[SYS_07_02]}
\reqverify{Испытание: Для «2 из 3» один голос оставляет ожидание, два закрывают запрос; начало этапа ещё проверяет предшественника и среду.}
\reqnote{Вход: голоса допустимых адресатов открытого запроса. Выход: «согласовано» при достижении N либо ожидание. Тип: функциональное. Статус: проект.}
\end{req}

\begin{req}{SRD_07_03}{1}
ПО должно принимать одно решение «согласовано» или «отклонено» от допустимого адресата открытого запроса, отклоняя повтор, постороннего адресата, неизвестное значение и закрытый запрос.
\reqsource{parent=[SYS_07_03]}
\reqverify{Испытание: Все перечисленные классы проверены; повтор не увеличивает число голосов.}
\reqnote{Вход: адресат, состояние запроса и значение решения. Выход: учтённое решение либо отказ. Тип: функциональное. Статус: проект.}
\end{req}

\begin{req}{SRD_07_04}{1}
ПО должно отклонять оба вида решения от инициатора продвижения релиза независимо от совмещения ролей.
\reqsource{parent=[SYS_07_04]}
\reqverify{Испытание: Инициатор с ролью согласующего пытается согласовать и отклонить; оба отказаны.}
\reqnote{Вход: решение инициатора продвижения. Выход: отказ обоим видам решения. Тип: функциональное. Статус: проект.}
\end{req}

\begin{req}{SRD_07_05}{1}
После первого допустимого отрицательного решения открытого запроса ПО должно закрывать его как «отклонено».
\reqsource{parent=[SYS_07_05]}
\reqverify{Испытание: При «2 из 3» один положительный и один отрицательный закрывают запрос; поздний положительный отказан.}
\reqnote{Вход: первое допустимое отрицательное решение открытого запроса. Выход: закрытый запрос «отклонено». Тип: функциональное. Статус: проект.}
\end{req}

\begin{req}{SRD_07_06}{1}
ПО должно принимать подтверждение выпуска только от деливери-менеджера после подтверждённой полной промышленной установки, успеха шагов и назначенной КТК, строго до сохранённого конца окна и до назначения отката.
\reqsource{parent=[SYS_07_06]}
\reqverify{Испытание: Полный набор принят; поочерёдное нарушение условий, включая момент ровно конца окна, даёт отказ.}
\reqnote{Вход: роль подтверждающего, результаты установки/шагов/КТК, срок окна и состояние отката. Выход: подтверждение выпуска либо отказ. Тип: функциональное. Статус: проект.}
\end{req}

\subsubsection{Учёт установок и откат}

\begin{req}{SRD_08_01}{1}
По запросу любой из восьми ролей ПО должно возвращать последние подтверждённые версии среды и время подтверждения, явно обозначая отсутствие подтверждения.
\reqsource{parent=[SYS_08_01]}
\reqverify{Испытание: После установки данные совпадают с ответом; неподтверждённая среда явно отмечена.}
\reqnote{Вход: запрос просмотра среды. Выход: последний подтверждённый состав, время и признак достоверности. Тип: функциональное. Статус: проект.}
\end{req}

\begin{req}{SRD_08_02}{1}
ПО должно обновлять подтверждённый состав только по достоверному результату установки, отката или подтверждению инженера эксплуатации.
\reqsource{parent=[SYS_08_02]}
\reqverify{Испытание: Успешные полные результаты меняют учёт; ошибка и потеря ответа сохраняют прежние сведения с пометкой неопределённости вместо выдачи цели за факт.}
\reqnote{Вход: результат установки, отката или проверки инженера. Выход: обновлённый подтверждённый состав либо прежние сведения с пометкой неопределённости. Тип: функциональное. Статус: проект.}
\end{req}

\begin{req}{SRD_08_03}{1}
По запросу релиз-менеджера при ошибке этапа после пройденного ПСИ ПО должно назначать ручной откат к сохранённому составу до установки, отклоняя неизвестную цель и ошибку до ПСИ или на ПСИ.
\reqsource{parent=[SYS_08_03]}
\reqverify{Испытание: После ПСИ с целью назначено; до ПСИ, на ПСИ и без цели отказ; без запроса операция отсутствует.}
\reqnote{Вход: запрос ручного отката, положение относительно ПСИ, выпуски, цель и допуск к среде. Выход: назначение/задание отката либо отказ/ожидание согласно карточке. Тип: функциональное. Статус: проект.}
\end{req}

\begin{req}{SRD_08_04}{1}
По запросу релиз-менеджера для последнего успешного промышленного выпуска с установленным в среде составом ПО должно назначать ручной откат к известному составу предшествующего успешного выпуска того же продукта.
\reqsource{parent=[SYS_08_03]}
\reqverify{Испытание: Для установленного R2 цель R1; старый выпуск, отсутствующий предшественник и несовпадающий состав отказаны.}
\reqnote{Вход: запрос ручного отката, положение относительно ПСИ, выпуски, цель и допуск к среде. Выход: назначение/задание отката либо отказ/ожидание согласно карточке. Тип: функциональное. Статус: проект.}
\end{req}

\begin{req}{SRD_08_05}{1}
ПО должно передавать ручное задание отката после завершения исходных заданий с известным результатом и выполнения условий доступа к среде.
\reqsource{parent=[SYS_08_03]}
\reqverify{Испытание: Неизвестный результат, чужая занятость и закрытое промышленное окно сохраняют ожидание; при допуске передана заданная цель.}
\reqnote{Вход: запрос ручного отката, положение относительно ПСИ, выпуски, цель и допуск к среде. Выход: назначение/задание отката либо отказ/ожидание согласно карточке. Тип: функциональное. Статус: проект.}
\end{req}

\begin{req}{SRD_08_06}{1}
После подтверждения восстановления всего целевого состава отката ПО должно переводить релиз из «Выпущен», «На продвижении» или «Требует вмешательства» в «Откачен».
\reqsource{parent=[SYS_08_04]}
\reqverify{Испытание: Каждое исходное состояние проверено для ручного и автоматического отката; частичный результат и отсутствие ответа перехода не дают.}
\reqnote{Вход: достоверный результат полного восстановления целевого состава. Выход: «Откачен». Тип: функциональное. Статус: проект.}
\end{req}

\begin{req}{SRD_08_07}{1}
ПО должно принимать ожидание отката как положительное число секунд: для ручного при запросе, для автоматического до запуска продвижения.
\reqsource{parent=[SYS_08_05]}
\reqverify{Испытание: Для обоих типов положительное значение принято, 0 и отрицательное отказаны.}
\reqnote{Вход: тип отката, ожидание в секундах и результат/истечение ожидания. Выход: принятое ожидание либо вмешательство/отказ согласно карточке. Тип: функциональное. Статус: проект.}
\end{req}

\begin{req}{SRD_08_08}{1}
При ошибке или истечении ожидания отката ПО должно обозначать операцию как требующую вмешательства, сохраняя прежнее состояние релиза при ручном откате и устанавливая «Требует вмешательства» при автоматическом.
\reqsource{parent=[SYS_08_05]}
\reqverify{Испытание: Оба типа проверить на ошибке и тайм-ауте от передачи задания; среда занята, «Откачен» отсутствует.}
\reqnote{Вход: тип отката, ожидание в секундах и результат/истечение ожидания. Выход: принятое ожидание либо вмешательство/отказ согласно карточке. Тип: функциональное. Статус: проект.}
\end{req}

\begin{req}{SRD_08_09}{1}
Перед первой установкой на этапе после ПСИ ПО должно сохранять подтверждённый исходный состав как цель отката без замены при повторах.
\reqsource{parent=[SYS_08_06]}
\reqverify{Испытание: После частичной установки и повтора цель прежняя; подтверждённая пустая цель допустима.}
\reqnote{Вход: подтверждённый состав перед первой установкой и последующие повторы. Выход: сохранённая неизменная цель либо блокировка промышленной установки. Тип: функциональное. Статус: проект.}
\end{req}

\begin{req}{SRD_08_10}{1}
При неизвестном исходном составе промышленной среды ПО должно блокировать начало установки.
\reqsource{parent=[SYS_08_06]}
\reqverify{Испытание: Без цели задания установки нет; подтверждённый пустой состав снимает именно это препятствие.}
\reqnote{Вход: подтверждённый состав перед первой установкой и последующие повторы. Выход: сохранённая неизменная цель либо блокировка промышленной установки. Тип: функциональное. Статус: проект.}
\end{req}

\begin{req}{SRD_08_11}{1}
При первой промышленной установке ПО должно сохранять конец использованного окна как крайний срок подтверждения без переноса при изменении расписания и повторе установки.
\reqsource{parent=[SYS_08_07]}
\reqverify{Испытание: Изменить окно и повторить установку; исходный срок сохранён.}
\reqnote{Вход: первая промышленная установка, конец окна, подтверждение и исходные задания. Выход: сохранённый срок, единственный автооткат либо ожидание его выдачи согласно карточке. Тип: функциональное. Статус: проект.}
\end{req}

\begin{req}{SRD_08_12}{1}
При достижении сохранённого конца окна без подтверждения выпуска и без завершённого восстановления ПО должно назначать единственный автооткат к составу до установки с прекращением новых промышленных шагов.
\reqsource{parent=[SYS_08_07]}
\reqverify{Испытание: Проверить «На продвижении» и «Требует вмешательства»; повтор события не дублирует откат, своевременное подтверждение и завершённое восстановление исключают назначение.}
\reqnote{Вход: первая промышленная установка, конец окна, подтверждение и исходные задания. Выход: сохранённый срок, единственный автооткат либо ожидание его выдачи согласно карточке. Тип: функциональное. Статус: проект.}
\end{req}

\begin{req}{SRD_08_13}{1}
ПО должно передавать назначенное автоматическое задание отката вне окна после завершения исходных заданий с известным результатом независимо от позднего подтверждения и изменения расписания.
\reqsource{parent=[SYS_08_07]}
\reqverify{Испытание: Установка завершена после окна: позднее подтверждение отказано, затем выдан один откат; неизвестный результат сохраняет ожидание выдачи.}
\reqnote{Вход: первая промышленная установка, конец окна, подтверждение и исходные задания. Выход: сохранённый срок, единственный автооткат либо ожидание его выдачи согласно карточке. Тип: функциональное. Статус: проект.}
\end{req}

\begin{req}{SRD_08_14}{1}
После ошибки отката ПО должно принимать разрешение дальнейших работ от инженера эксплуатации только после завершения всех внешних заданий среды с известным результатом. Для разрешения указываются фактические версии, время проверки и основание подтверждения.
\reqsource{parent=[SYS_08_08]}
\reqverify{Испытание: Полные данные снимают блокировку и подтверждают состав; неизвестное задание, неполные данные и другая роль отказаны; «Откачен» этим не присваивается.}
\reqnote{Вход: версии, время и основание проверки инженера, конечные результаты заданий. Выход: подтверждённый состав и разрешение работ либо отказ. Тип: функциональное. Статус: проект.}
\end{req}

\subsubsection{Управление доступом}

\begin{req}{SRD_09_01}{1}
Перед пользовательской операцией ПО должно проверять подлинность, срок действия Bearer-токена и соответствие подтверждённой личности пользователю запроса.
\reqsource{parent=[SYS_09_01]}
\reqverify{Испытание: Действительный токен допускает проверку ролей; чужая личность отказана.}
\reqnote{Вход: пользовательский запрос, Bearer-токен и результат проверки службой. Выход: подтверждённая личность для проверки прав либо отказ. Тип: функциональное. Статус: проект.}
\end{req}

\begin{req}{SRD_09_02}{1}
При отсутствии, подделке, истечении токена или отсутствии достоверного подтверждения его действительности ПО должно отклонять операцию.
\reqsource{parent=[SYS_09_01]}
\reqverify{Испытание: Каждый класс проверен; предметных изменений нет.}
\reqnote{Вход: пользовательский запрос, Bearer-токен и результат проверки службой. Выход: подтверждённая личность для проверки прав либо отказ. Тип: функциональное. Статус: проект.}
\end{req}

\begin{req}{SRD_09_03}{1}
По запросу администратора ПО должно назначать одну или несколько из восьми ролей пользователю, отклоняя неизвестную роль.
\reqsource{parent=[SYS_09_02]}
\reqverify{Испытание: Две роли сохранены; неизвестная роль отказана.}
\reqnote{Вход: администратор, пользователь и назначаемые роли. Выход: набор ролей либо отказ неизвестному значению. Тип: функциональное. Статус: проект.}
\end{req}

\begin{req}{SRD_09_04}{1}
ПО должно допускать операцию по объединению ролей таблицы пользователей с сохранением ограничений исполнителя и инициатора продвижения.
\reqsource{parent=[SYS_09_03]}
\reqverify{Испытание: Аудитор не запускает релиз; администратор без роли аудитора не читает журнал; совмещение сохраняет запрет самосогласования; проверить всю матрицу ролей.}
\reqnote{Вход: личность, роли, действие, назначенный исполнитель и инициатор продвижения. Выход: допуск либо отказ в полномочиях. Тип: функциональное. Статус: проект.}
\end{req}

\subsubsection{Журнал аудита}

\begin{req}{SRD_10_01}{1}
При каждом изменении предметных данных ПО должно добавлять запись аудита с датой, временем, инициатором, объектом, действием и результатом.
\reqsource{parent=[SYS_10_01]}
\reqverify{Испытание: Изменить каталог, роли, состав, маршрут, решение и установку; все поля отражены; автоматический инициатор --- система с событием; запись аудита не порождает запись о себе.}
\reqnote{Вход: предметное изменение с инициатором и результатом. Выход: запись даты, времени, инициатора, объекта, действия и результата. Тип: функциональное. Статус: проект.}
\end{req}

\begin{req}{SRD_10_02}{1}
ПО должно отклонять редактирование и удаление записей аудита через все предоставляемые им операции независимо от роли.
\reqsource{parent=[SYS_10_02]}
\reqverify{Испытание: Каждая роль, включая администратора, получает отказ; исходные записи сохранены.}
\reqnote{Вход: запрос изменения или удаления записи аудита от любой роли. Выход: отказ; исходная запись сохранена. Тип: функциональное. Статус: проект.}
\end{req}

\begin{req}{SRD_10_03}{1}
По запросу аудитора ПО должно возвращать записи аудита в хронологическом порядке с сохранёнными временем и инициаторами.
\reqsource{parent=[SYS_10_03]}
\reqverify{Испытание: Изменения разных моментов возвращены по времени без изменения значений.}
\reqnote{Вход: запрос просмотра аудитора. Выход: журнал в хронологическом порядке. Тип: функциональное. Статус: проект.}
\end{req}

\subsection{Требования к производительности}
Время измеряется на границе ПО. Имитаторы доступны; исходные данные и последовательность нагрузки сохраняются с протоколом для повторения прогона. Числовые пороги взяты из системных требований, а не назначены заново.

\begin{req}{SRD_11_01}{1}
ПО должно делать зарегистрированную версию доступной не позднее 60 с от получения полных допустимых сведений о сборке при продвижении до 20 релизов одновременно.
\reqsource{parent=[SYS_11_01]}
\reqverify{Испытание: При указанной нагрузке каждый интервал от приёма до доступности в каталоге не больше 60 с.}
\reqnote{Вход: полное сообщение сборки при продвижении до 20 релизов. Выход: версия доступна в каталоге не позднее 60 с. Тип: характеристика качества. Статус: проект.}
\end{req}

\begin{req}{SRD_11_02}{1}
ПО должно поддерживать одновременное продвижение 20 релизов с независимыми целевыми средами и доступными внешними системами.
\reqsource{parent=[SYS_11_02]}
\reqverify{Испытание: Все 20 проходят маршруты без отказа из-за числа релизов.}
\reqnote{Вход: 20 релизов с независимыми средами и доступными внешними системами. Выход: одновременное продвижение всех 20 релизов. Тип: характеристика качества. Статус: проект.}
\end{req}

\begin{req}{SRD_11_03}{1}
ПО должно возвращать просмотр каталога, состава, состояния релиза либо аудита не позднее 3 с от приёма запроса при 30 работающих пользователях и 20 продвигаемых релизах.
\reqsource{parent=[SYS_11_04]}
\reqverify{Испытание: 30 одновременных запросов при 20 продвижениях включают все четыре вида просмотра; каждый ответ за 3 с или меньше без сетевой задержки до пользователя.}
\reqnote{Вход: 30 запросов пользователей при продвижении 20 релизов. Выход: каждый результат просмотра не позднее 3 с. Тип: характеристика качества. Статус: проект.}
\end{req}

\subsection{Атрибуты качества}
\subsubsection{Надёжность и восстановление}

\begin{req}{SRD_11_04}{1}
После перезапуска при сохранном хранилище ПО должно восстанавливать до возобновления продвижения все подтверждённые состояния релизов, результаты шагов, решения, подтверждения выпуска, сроки окон, цели и операции отката и аудит.
\reqsource{parent=[SYS_11_03]}
\reqverify{Испытание: Перезапустить на шаге, согласовании и после фиксации результата; утрачено 0 подтверждённых записей.}
\reqnote{Вход: сохранные данные при перезапуске, незавершённые задания и сроки окон. Выход: восстановленное состояние, запрос состояния задания либо назначение автоотката согласно карточке. Тип: характеристика качества. Статус: проект.}
\end{req}

\begin{req}{SRD_11_05}{1}
После восстановления ПО должно сохранять ожидание каждого незавершённого внешнего задания до установления его состояния запросом исходного задания.
\reqsource{parent=[SYS_11_03,SYS_04_10]}
\reqverify{Испытание: Остановить после передачи установки с потерей ответа; повторных запусков до выяснения результата 0.}
\reqnote{Вход: сохранные данные при перезапуске, незавершённые задания и сроки окон. Выход: восстановленное состояние, запрос состояния задания либо назначение автоотката согласно карточке. Тип: характеристика качества. Статус: проект.}
\end{req}

\begin{req}{SRD_11_06}{1}
Если промышленное окно истекло во время остановки, ПО должно после восстановления назначать автооткат при отсутствии своевременного подтверждения и завершённого восстановления.
\reqsource{parent=[SYS_11_03,SYS_08_07]}
\reqverify{Испытание: Остановить перед концом окна, запустить после; автооткат назначен, неизвестное исходное задание блокирует выдачу отката.}
\reqnote{Вход: сохранные данные при перезапуске, незавершённые задания и сроки окон. Выход: восстановленное состояние, запрос состояния задания либо назначение автоотката согласно карточке. Тип: характеристика качества. Статус: проект.}
\end{req}

\subsubsection{Безопасность и сопровождаемость}
Проверку доступа задают карточки \reqref{SRD_09_01}, \reqref{SRD_09_02}, \reqref{SRD_09_03} и \reqref{SRD_09_04}. Запрет самосогласования задан в \reqref{SRD_07_04}, защита аудита --- в \reqref{SRD_10_02}. Защита от разрушения хранилища за границами версии 1. Защита сетевого обмена требует согласования по разделу 3.1.4. Сопровождение требований обеспечивают постоянные ID, версии и матрица~\ref{tab:trace}; самостоятельного числового показателя сопровождаемости продукта во входных требованиях нет.
\subsection{Матрица трассировки}
Таблица~\ref{tab:trace} содержит все системные требования версии 4 и уточняющие карточки СПТ. Связь не доказывает реализацию либо прохождение испытаний. Производных требований в этой редакции нет. Контракты связи остаются открытым вопросом; до их согласования интерфейсный раздел не считается завершённым.
\begin{tabl}{|P{3.3cm}|L|}{Системные требования и требования к ПО}[tab:trace]{\textbf{Системное требование} & \textbf{Требования к ПО}}

\reqref{SYS_01_01} & \reqref{SRD_01_01} \\ \hline

\reqref{SYS_01_02} & \reqref{SRD_01_02} \\ \hline

\reqref{SYS_01_03} & \reqref{SRD_01_03} \\ \hline

\reqref{SYS_01_04} & \reqref{SRD_01_04} \\ \hline

\reqref{SYS_01_05} & \reqref{SRD_01_05}, \reqref{SRD_01_06} \\ \hline

\reqref{SYS_01_06} & \reqref{SRD_01_07}, \reqref{SRD_01_08} \\ \hline

\reqref{SYS_01_07} & \reqref{SRD_01_09} \\ \hline

\reqref{SYS_02_01} & \reqref{SRD_02_01} \\ \hline

\reqref{SYS_02_02} & \reqref{SRD_02_02} \\ \hline

\reqref{SYS_02_03} & \reqref{SRD_02_03} \\ \hline

\reqref{SYS_02_04} & \reqref{SRD_02_04}, \reqref{SRD_02_05} \\ \hline

\reqref{SYS_02_05} & \reqref{SRD_02_06} \\ \hline

\reqref{SYS_02_06} & \reqref{SRD_02_07} \\ \hline

\reqref{SYS_02_07} & \reqref{SRD_02_08} \\ \hline

\reqref{SYS_02_08} & \reqref{SRD_02_09} \\ \hline

\reqref{SYS_03_01} & \reqref{SRD_03_01} \\ \hline

\reqref{SYS_03_02} & \reqref{SRD_03_02} \\ \hline

\reqref{SYS_03_03} & \reqref{SRD_03_03} \\ \hline

\reqref{SYS_03_04} & \reqref{SRD_03_04} \\ \hline

\reqref{SYS_03_05} & \reqref{SRD_03_05} \\ \hline

\reqref{SYS_03_06} & \reqref{SRD_03_06} \\ \hline

\reqref{SYS_04_01} & \reqref{SRD_04_01} \\ \hline

\reqref{SYS_04_02} & \reqref{SRD_04_02} \\ \hline

\reqref{SYS_04_03} & \reqref{SRD_04_03} \\ \hline

\reqref{SYS_04_04} & \reqref{SRD_04_04}, \reqref{SRD_04_06} \\ \hline

\reqref{SYS_04_05} & \reqref{SRD_04_05}, \reqref{SRD_04_06}, \reqref{SRD_04_08} \\ \hline

\reqref{SYS_04_06} & \reqref{SRD_04_09} \\ \hline

\reqref{SYS_04_07} & \reqref{SRD_04_10}, \reqref{SRD_04_11} \\ \hline

\reqref{SYS_04_08} & \reqref{SRD_04_12} \\ \hline

\reqref{SYS_04_09} & \reqref{SRD_04_13}, \reqref{SRD_04_14} \\ \hline

\reqref{SYS_04_10} & \reqref{SRD_04_06}, \reqref{SRD_04_07}, \reqref{SRD_04_08}, \reqref{SRD_11_05} \\ \hline

\reqref{SYS_04_11} & \reqref{SRD_04_15} \\ \hline

\reqref{SYS_04_12} & \reqref{SRD_04_16} \\ \hline

\reqref{SYS_04_13} & \reqref{SRD_04_17} \\ \hline

\reqref{SYS_05_01} & \reqref{SRD_05_01}, \reqref{SRD_05_02}, \reqref{SRD_05_03}, \reqref{SRD_05_04} \\ \hline

\reqref{SYS_05_02} & \reqref{SRD_05_05} \\ \hline

\reqref{SYS_05_03} & \reqref{SRD_05_06}, \reqref{SRD_05_07} \\ \hline

\reqref{SYS_06_01} & \reqref{SRD_06_01} \\ \hline

\reqref{SYS_06_02} & \reqref{SRD_06_02} \\ \hline

\reqref{SYS_06_03} & \reqref{SRD_06_03}, \reqref{SRD_06_04}, \reqref{SRD_06_05} \\ \hline

\reqref{SYS_06_04} & \reqref{SRD_06_06} \\ \hline

\reqref{SYS_07_01} & \reqref{SRD_07_01} \\ \hline

\reqref{SYS_07_02} & \reqref{SRD_07_02} \\ \hline

\reqref{SYS_07_03} & \reqref{SRD_07_03} \\ \hline

\reqref{SYS_07_04} & \reqref{SRD_07_04} \\ \hline

\reqref{SYS_07_05} & \reqref{SRD_07_05} \\ \hline

\reqref{SYS_07_06} & \reqref{SRD_07_06} \\ \hline

\reqref{SYS_08_01} & \reqref{SRD_08_01} \\ \hline

\reqref{SYS_08_02} & \reqref{SRD_08_02} \\ \hline

\reqref{SYS_08_03} & \reqref{SRD_08_03}, \reqref{SRD_08_04}, \reqref{SRD_08_05} \\ \hline

\reqref{SYS_08_04} & \reqref{SRD_08_06} \\ \hline

\reqref{SYS_08_05} & \reqref{SRD_08_07}, \reqref{SRD_08_08} \\ \hline

\reqref{SYS_08_06} & \reqref{SRD_08_09}, \reqref{SRD_08_10} \\ \hline

\reqref{SYS_08_07} & \reqref{SRD_08_11}, \reqref{SRD_08_12}, \reqref{SRD_08_13}, \reqref{SRD_11_06} \\ \hline

\reqref{SYS_08_08} & \reqref{SRD_08_14} \\ \hline

\reqref{SYS_09_01} & \reqref{SRD_09_01}, \reqref{SRD_09_02} \\ \hline

\reqref{SYS_09_02} & \reqref{SRD_09_03} \\ \hline

\reqref{SYS_09_03} & \reqref{SRD_09_04} \\ \hline

\reqref{SYS_10_01} & \reqref{SRD_10_01} \\ \hline

\reqref{SYS_10_02} & \reqref{SRD_10_02} \\ \hline

\reqref{SYS_10_03} & \reqref{SRD_10_03} \\ \hline

\reqref{SYS_11_01} & \reqref{SRD_11_01} \\ \hline

\reqref{SYS_11_02} & \reqref{SRD_11_02} \\ \hline

\reqref{SYS_11_03} & \reqref{SRD_11_04}, \reqref{SRD_11_05}, \reqref{SRD_11_06} \\ \hline

\reqref{SYS_11_04} & \reqref{SRD_11_03} \\ \hline

\end{tabl}
\section{Приложение: иллюстрации}
Таблица~\ref{tab:scenarios} поясняет совместное действие карточек, не вводит новых требований и не заменяет будущие тест-кейсы.
\begin{tabl}{|P{3.6cm}|L|}{Сценарии и граничные случаи}[tab:scenarios]{\textbf{Сценарий} & \textbf{Последовательность}}
Штатный выпуск & Фиксация, запуск, согласование, свободная среда, шаги и КТК. Ручная приёмка ПСИ. Промышленная установка в окне; подтверждение деливери-менеджера до конца; «Выпущен». \\ \hline
Потеря ответа & Результат неизвестен, следующий шаг и повтор заблокированы; запрос состояния исходного задания. Перезапуск сохраняет ту же идентичность. \\ \hline
Повтор после ошибки & Известен результат, среда доступна, согласование действует, отката нет. Новая попытка прежнего шага; следующие результаты только в истории; КТК вычисляется заново. \\ \hline
Конец окна & Начатое задание не прерывается. Без подтверждения назначен автооткат; его выдача ждёт известного завершения исходного задания. Затем один откат вне окна. \\ \hline
Ошибка после ПСИ & Только запрос релиз-менеджера назначает ручной откат к исходному составу. В промышленной среде начало ожидает окна. \\ \hline
Ошибка отката & Среда занята; после известных конечных результатов инженер эксплуатации подтверждает фактический состав и разрешает работы. Это не переводит релиз в «Откачен». \\ \hline
Отказ согласующего & Запрос закрыт с отказом; поздние голоса отклонены. Повтор шага не обходит отклонённое согласование. \\ \hline
\end{tabl}
\begin{changelog}
\change{1}{13.11.2026}{Первая версия СПТ: входы и выходы, штатные и ошибочные режимы, логический обмен, показатели качества и обратная трассировка. Сетевые контракты вынесены на согласование}{Карцева П. А.}{Задание к 14.11.2026; системные требования v4 и стандарт требований v2}
\end{changelog}
\end{document}
